\documentclass[../zusammenfassung_AN.tex]{subfiles}
\begin{document}
	
	\subsection{Potenzreihenentwicklung}
	
	\subsubsection{Faltung (Cauchy-Produkt) von Potenzreihen}
	
	\begin{tcolorbox}[dglbox, title={Cauchy-Produkt}]
		Gegeben zwei Potenzreihen
		\[
		A(x) = \sum_{n=0}^{\infty} a_n x^n, \qquad B(x) = \sum_{n=0}^{\infty} b_n x^n,
		\]
		ist ihr Produkt wieder eine Potenzreihe:
		\[
		A(x)\cdot B(x) = \sum_{n=0}^{\infty} c_n x^n, \qquad
		\boxed{c_n = \sum_{k=0}^{n} a_k \cdot b_{n-k}}
		\]
	\end{tcolorbox}
	
	\paragraph{Faltungsschema} — $c_n$ ist die Summe entlang der Gegendiagonale $k+j=n$:
	
	\[
	\begin{array}{c|ccccc}
		\cdot  & b_0       & b_1       & b_2       & b_3       & \cdots \\\hline
		a_0    & a_0 b_0   & a_0 b_1   & a_0 b_2   & a_0 b_3   & \cdots \\
		a_1    & a_1 b_0   & a_1 b_1   & a_1 b_2   & \cdots    &        \\
		a_2    & a_2 b_0   & a_2 b_1   & \cdots    &           &        \\
		a_3    & a_3 b_0   & \cdots    &           &           &        \\
		\vdots & \vdots    &           &           &           &
	\end{array}
	\]
	
	\begin{tcolorbox}[formelbox]
		\textbf{Erste Koeffizienten explizit:}
		\begin{align*}
			c_0 &= a_0 b_0 \\
			c_1 &= a_0 b_1 + a_1 b_0 \\
			c_2 &= a_0 b_2 + a_1 b_1 + a_2 b_0 \\
			c_3 &= a_0 b_3 + a_1 b_2 + a_2 b_1 + a_3 b_0
		\end{align*}
		\textbf{Merkregel:} In jedem Term von $c_n$ addieren sich die Indizes zu $n$:
		\[
		c_n = \sum_{\substack{k+j=n \\ k,j\,\geq\,0}} a_k\,b_j
		\]
	\end{tcolorbox}
	
	\begin{tcolorbox}[hinweis]
		\textbf{Gültigkeitsbereich:} Das Cauchy-Produkt konvergiert gegen $A(x)\cdot B(x)$,
		wenn mindestens eine der beiden Reihen absolut konvergiert (Satz von Mertens).\\
		Sicher anwendbar für $|x-x_0| < \min(R_A,\,R_B)$.
	\end{tcolorbox}
	
\end{document}